\begin{codebox}
\codeline{\textcolor{cmtgray}{\#\ ==============================}}
\codeline{\textcolor{cmtgray}{\#\ 2.\ 工程本体论的转向}}
\codeline{\textcolor{cmtgray}{\#\ ==============================}}
\codeline{\textcolor{cmtgray}{\#\ 知识工程领域（小写复数的\ ontologies）完成了一次实用主义转向：}}
\codeline{\textcolor{cmtgray}{\#\ 不再追问存在的本质，而是约定——}}
\codeblank{}
\codeline{\textcolor{cmtgray}{\#\ \ \ 在一个领域内，存在哪些概念（类）？}}
\codeline{\textcolor{cmtgray}{\#\ \ \ 概念之间有什么关系（属性）？}}
\codeline{\textcolor{cmtgray}{\#\ \ \ 什么样的陈述是合法的（公理与约束）？}}
\codeblank{}
\codeline{\textcolor{cmtgray}{\#\ Gruber\ (1993)\ 的经典定义：}}
\codeline{\textcolor{cmtgray}{\#\ \ \ "An\ ontology\ is\ an\ explicit\ specification\ of\ a\ conceptualization."}}
\codeline{\textcolor{cmtgray}{\#\ \ \ 本体是概念化的显式规范}}
\codeblank{}
\codeline{\textcolor{cmtgray}{\#\ Studer\ et\ al.\ (1998)\ 的综合定义包含四个关键词：}}
\codeline{\textcolor{cmtgray}{\#\ \ \ 概念化（Conceptualization）——\ 对世界某侧面的抽象模型}}
\codeline{\textcolor{cmtgray}{\#\ \ \ 显式（Explicit）\ \ \ \ \ \ \ \ \ \ ——\ 概念与约束被明确定义，而非隐含}}
\codeline{\textcolor{cmtgray}{\#\ \ \ 形式化（Formal）\ \ \ \ \ \ \ \ \ \ ——\ 机器可读、可推理，而非自然语言描述}}
\codeline{\textcolor{cmtgray}{\#\ \ \ 共享（Shared）\ \ \ \ \ \ \ \ \ \ \ \ ——\ 团体共识，而非个人观点}}
\codeblank{}
\end{codebox}
